\scanpage{75}
\mspage{74}

Conditions
\begin{tabular}{|p{0.92\linewidth}}
\[
\text{Soient }\unsure{x|a},y\in L,\quad\text{alors}\quad
x\parallel y\quad\Longleftrightarrow\quad |x|\cap|y|=\varnothing
\tag*{$(\mathrm{i})_L$}
\]
\[
\begin{aligned}
&\deleted{\text{\unsure{Comme}}}\quad\forall\unsure{\xi|\zeta},\eta\in P,
\ \inserted{\xi\ne\eta},\ \exists\deleted{\unsure{y}}x\in L,\quad\text{avec}\quad
\xi\in|x|,\ \eta\notin|x|\\
&\hspace{27em}\text{ou l'inverse}
\end{aligned}
\tag*{$(\mathrm{ii})_L$}
\]

\end{tabular}

\uline{Exemple}. $P$ est une ficelle topologique, \deleted{\unsure{orientée}} et $M$ est le \unsure{stratos} des \inserted{liens et des} segments (ou sous-ficelles) de ladite, la fonction support $|x|$ \deleted{est} pour les liens étant la fonction évidente.
% slanted
\margin{\uline{NB} Il faut supposer ici pour la validité de (i), que les ``bouts'' de la ficelle soient connexes}

La condition (ii) n'est pas forcément satisfaite (sûrement pas, si les bouts de la ficelle ne \unsure{sont|st} \unsure{réduits} à des points --- s'ils le sont, on a une bonne chance pour (ii) soit valable\ldots)

Dans tous les cas intéressants des \unsure{stratos|Atlas} à ma connaissance, sauf peut être celui \unsure{de l'Atlas|des Atlas} associé à un ens. ordonné (\unsure{toujours} si on veut) donné ad hoc, on a une notion naturelle de supports $|\ |$.

La question qui se pose alors \unsure{est ici|est bien}: Une telle donnée détermine-t-elle $P$ à isomorphisme canonique près --- p. ex. par \unsure{l'existence} d'un iso can. d'ens. ordonnés ou \deleted{application} \inserted{\unsure{anti-isomorphisme}} \unsure{(}
\[
\underset{\text{ens des supports ``abstraits'' de }M}{\Sigma_M}
\simeq\unsure{\mathcal{P}|\mathscr{P}}(P)\quad??
\]
